\chapter{Sécurité des échanges}\label{chap:securite}

\section{Authentification}
\begin{description}[style=nextline]
  \item[Services SOAP — WS-Security] Chaque requête doit porter un en-tête OASIS WS-Security 1.0
        \code{<wsse:Security><wsse:UsernameToken>} déclaré \code{mustUnderstand}. Un intercepteur JAX-WS
        (\code{SoapSecurityHandler}) extrait le jeton, authentifie l'appelant et dépose l'utilisateur dans le contexte du
        message. En cas d'échec, la requête n'atteint pas le service : la réponse est une SOAP Fault
        \code{wsse:FailedAuthentication}. Le WSDL reste consultable sans authentification.
  \item[Services REST — clé d'API ou HTTP Basic] Un filtre JAX-RS (\code{AuthentificationFilter}, priorité
        \code{AUTHENTICATION}) accepte deux schémas : une \textbf{clé d'API} dans l'en-tête \code{X-API-Key}, destinée aux
        applications (orchestrateur, Postman, scripts), ou HTTP Basic, utilisé par l'interface web après connexion.
        Une clé invalide ou révoquée est refusée (\code{401 CLE\_API\_INVALIDE}). Seul \code{/api/sante} est public.
  \item[Gestion des clés] Seule l'empreinte SHA-256 de la clé est stockée. \code{POST /api/moi/cle-api} (bouton
        « Ma clé d'API » de l'interface), réservé au responsable des stages (403 pour les autres rôles), génère une
        nouvelle clé aléatoire de 256 bits, l'affiche une seule fois et révoque la précédente.
  \item[Appels inter-services] L'orchestrateur et les services composites utilisent un compte technique de rôle
        \code{SYSTEME}, configurable, injecté par un intercepteur client (SOAP) ou par sa clé d'API (REST).
\end{description}

\section{Protection des informations sensibles}
\begin{itemize}
  \item \textbf{Mots de passe} hachés par PBKDF2-HMAC-SHA256 (120\,000 itérations, sel aléatoire de 128 bits) ;
        la comparaison s'effectue en temps constant (\code{MessageDigest.isEqual}). Seule une empreinte SHA-256 du couple
        identifiant/mot de passe est gardée en cache mémoire pour éviter de refaire la dérivation à chaque appel.
  \item \textbf{Masquage des données personnelles} : lorsqu'une entreprise consulte la fiche d'un candidat
        (\code{consulterEtudiant}), elle voit son profil académique (filière, niveau, compétences) mais le téléphone et
        l'e-mail sont masqués (\code{+212 6 \textbullet\textbullet{} \textbullet\textbullet{} \textbullet\textbullet{} 01}, \code{y\textbullet\textbullet\textbullet{}@…}) et l'adresse n'est pas transmise. Les
        adresses e-mail sont également masquées dans les journaux du serveur.
  \item \textbf{Cloisonnement des données} : un étudiant ne peut lire que sa fiche, ses candidatures, ses notifications et
        ses stages ; une entreprise que ses offres, les candidatures et stages qui la concernent, et elle ne peut pas modifier
        son propre référencement. Dans « Mon profil », chacun ne voit et ne modifie que sa propre fiche : l'étudiant ses
        coordonnées et compétences (identité, filière, niveau et statut restent gérés par l'école), l'entreprise ses
        coordonnées et son contact (raison sociale et secteur en lecture seule) ; le contrôle est fait par les services SOAP. Les filtres de recherche sont \emph{imposés} côté serveur, indépendamment des paramètres reçus.
  \item \textbf{Variables de processus} : l'API d'orchestration n'expose qu'une liste blanche de variables.
  \item \textbf{Requêtes SQL paramétrées} partout (\code{PreparedStatement}), ce qui exclut l'injection SQL.
  \item \textbf{Interface web} : toutes les données affichées sont échappées (prévention XSS), l'en-tête
        \code{X-Content-Type-Options: nosniff} est envoyé et les identifiants ne sont conservés qu'en \code{sessionStorage}.
\end{itemize}

\section{Contrôle d'accès par rôle}
\begin{table}[htbp]
\centering\small
\caption{Matrice des droits (L : lecture, E : écriture, P : périmètre propre uniquement)}
\begin{tabularx}{\textwidth}{@{}Ycccc@{}}
\toprule
\textbf{Ressource / opération} & \textbf{Étudiant} & \textbf{Entreprise} & \textbf{Responsable} & \textbf{Système} \\
\midrule
Fiche étudiant (consulter) & L (P) & L & L/E & L/E statut \\
Entreprise (fiche) & L & L/E (P, hors référencement) & L/E & L \\
Offres (créer, modifier, supprimer) & L & L/E (P) & L/E & L, réserver \\
Soumettre une candidature (processus) & E (P) & — & E & — \\
Candidatures (consulter) & L (P) & L (P) & L & L/E \\
Annuler / compléter une candidature & E (P) & — & E & E \\
Décision sur une candidature & — & E (P) & E & — \\
Notifications & L (P) & L (P) & L/E & E \\
Stages, comptes rendus & L/E (P) & L/E (P) & L/E & E \\
Évaluations & L (P) & E (P) & E & — \\
Incidents, instances, pannes simulées & — & — & L/E & — \\
\bottomrule
\end{tabularx}
\end{table}

\section{Confidentialité du transport}
UsernameToken en \code{PasswordText} et HTTP Basic transmettent le mot de passe encodé mais non chiffré : en
production, la plateforme doit être exposée exclusivement en \textbf{HTTPS} (terminaison TLS par un proxy inverse ou
\code{HttpsServer} du JDK), ce qui ne demande aucune modification des services. Le contexte de sécurité REST signale
d'ailleurs si la requête est sécurisée (\code{isSecure()}). Des évolutions possibles sont présentées au
chapitre~\ref{chap:conclusion} (jetons JWT/OAuth 2, signature WS-Security).
